\documentclass{article}
\usepackage[most]{tcolorbox}   % for the framed style, used below
\usepackage[parsep=0em,itemsep=0em]{propositions}

\begin{document}

\section{Test: \texttt{continue} key}

\subsection{Two consecutive prop environments (default: continue=true)}

Before text.

\begin{prop}
  \item First prop, item A.
  \item First prop, item B.
\end{prop}

\begin{prop}
  \item Second prop, item C.
  \item Second prop, item D.
\end{prop}

After text.

\subsection{Three consecutive prop environments}

Before text.

\begin{prop}
  \item Prop 1, item A.
\end{prop}

\begin{prop}
  \item Prop 2, item B.
\end{prop}

\begin{prop}
  \item Prop 3, item C.
\end{prop}

After text.

\subsection{Two consecutive prop environments with blank line between (should still continue)}

Before text.

\begin{prop}
  \item First prop, item A.
  \item First prop, item B.
\end{prop}

\begin{prop}
  \item Second prop, item C.
  \item Second prop, item D.
\end{prop}

After text.

\subsection{Intervening text breaks the chain}

Before text.

\begin{prop}
  \item First prop, item A.
\end{prop}

Some intervening text here.

\begin{prop}
  \item Second prop — should have normal topsep spacing above.
\end{prop}

After text.

\subsection{With continue=false globally}

\propoptions{continue=false}

Before text.

\begin{prop}
  \item First prop, item A.
\end{prop}

\begin{prop}
  \item Second prop — should have normal topsep spacing (continue=false).
\end{prop}

After text.

\propoptions{continue=true}

\subsection{Per-environment continue=false}

Before text.

\begin{prop}
  \item First prop, item A.
\end{prop}

\begin{prop}[continue=false]
  \item Second prop — should have normal topsep spacing.
\end{prop}

\begin{prop}
  \item Third prop — continue=true again (but previous broke chain, should be normal).
\end{prop}

After text.

\subsection{Nested prop environments}

Before text.

\begin{prop}
  \item Outer item 1.
  \begin{prop}
    \item Inner item A.
  \end{prop}
  \begin{prop}
    \item Inner item B — should continue from inner A.
  \end{prop}
  \item Outer item 2.
\end{prop}

After text.

\subsection{A \texttt{framed} list is not joined}

The built-in \texttt{framed} style carries \texttt{continue=false}, so a frame
is kept clear of any list immediately before or after it: joining would cancel
the \textsf{tcolorbox}'s own spacing rather than merely retune the gap.  The
first and third lists below must each be separated from the frame; the last two
must still join each other, since \texttt{continue=false} belongs to the framed
environment alone.

Before text.

\begin{prop}\item Before the frame.\end{prop}
\begin{prop}[style=framed]\item A framed list.\end{prop}
\begin{prop}\item After the frame.\end{prop}
\begin{prop}\item Joins the list above.\end{prop}

After text.

\subsection{Reference: single prop environment}

Before text.

\begin{prop}
  \item Item A.
  \item Item B.
  \item Item C.
\end{prop}

After text.

\subsection{A document \texttt{\textbackslash parskip} does not widen the join}

The gap a join imitates is the one between two items of a single environment,
and it must stay that gap however wide the surrounding paragraph spacing is.
The three pairs below set \verb|\parskip| to 0pt, 6pt and 24pt.  In each pair
the gap between B and C, which crosses an environment boundary, must match the
gap between A and B, which does not --- and the two gaps must look the same in
all three pairs, since neither is made of \verb|\parskip|.

\begingroup\setlength{\parskip}{0pt}
\begin{prop}\item Item A, \verb|\parskip=0pt|.\item Item B.\end{prop}
\begin{prop}\item Item C: level with B as B is with A.\end{prop}
\endgroup

\begingroup\setlength{\parskip}{6pt}
\begin{prop}\item Item A, \verb|\parskip=6pt|.\item Item B.\end{prop}
\begin{prop}\item Item C: level with B as B is with A.\end{prop}
\endgroup

\begingroup\setlength{\parskip}{24pt}
\begin{prop}\item Item A, \verb|\parskip=24pt|.\item Item B.\end{prop}
\begin{prop}\item Item C: level with B as B is with A.\end{prop}
\endgroup

The same inside an enclosing list, where the \verb|\parskip| the join has to
keep out is the outer list's \verb|parsep| rather than the document's.

\begin{prop}
  \item An outer item containing:
  \begin{prop}\item Item A.\item Item B.\end{prop}
  \begin{prop}\item Item C: level with B as B is with A.\end{prop}
  \item A further outer item.
\end{prop}

\end{document}
